% Part: applied-modal-logics
% Chapter: temporal-logic
% Section: introduction

\documentclass[../../../include/open-logic-section]{subfiles}

\begin{document}

\olfileid[hi]{aml}{tl}{int}

\olsection{परिचय}

कालिक तर्कशास्त्र उन दावों से संबंधित हैं जो बताते हैं कि कोई बात भविष्य
में होगी या अतीत में रही है। आर्थर प्रायर को कालिक तर्कशास्त्र का प्रवर्तक
माना जाता है; उन्होंने इसे \emph{काल-तर्कशास्त्र} कहा। यहाँ हमारा विवेचन
मुख्यतः प्रायर के मूल मोडल दृष्टिकोण का अनुसरण करेगा, जिसमें कालिक संकारकों
को प्रतिज्ञप्तिक तर्कशास्त्र के आधारभूत ढाँचे में जोड़ा जाता है; प्रतिज्ञप्तिक
तर्कशास्त्र दावों को सामान्यतः काल-रहित मानता है।

उदाहरण के लिए, प्रतिज्ञप्तिक तर्कशास्त्र में मैं बीज़ी नामक एक कुत्ते की बात
कर सकता हूँ, जो कभी बैठता है और कभी नहीं बैठता---जैसा कि कुत्ते प्रायः करते
हैं। शास्त्रीय तर्कशास्त्र में यह दावा करना विरोधाभासी होगा कि बीज़ी बैठा है
और यह भी कि बीज़ी बैठा नहीं है। किंतु स्पष्टतः दोनों बातें सत्य हो सकती हैं,
बस एक ही समय पर नहीं; भाषा में कालिक संकारक जोड़ने से हम इस दावे को अपेक्षाकृत
आसानी से व्यक्त कर सकते हैं। कालिक संकारकों को जोड़ने से हम ऐसे अनुमान की
मान्यता भी समझा सकते हैं जो ``बीज़ी को खाने का इनाम या गेंद मिलेगी'' से
``बीज़ी को खाने का इनाम मिलेगा या बीज़ी को गेंद मिलेगी'' तक जाता है।

फिर भी, कालिक तर्कशास्त्र अनेक दार्शनिक प्रश्न उठाता है, जिनके कारण हम उसके
एक ढाँचे को दूसरे ढाँचे के स्थान पर अपना सकते हैं। उदाहरण के लिए,
\emph{भावी आकस्मिक कथन} भविष्य के बारे में ऐसा कथन है जो न तो आवश्यक है और
न असंभव। यदि हम कहते हैं, ``रिचर्ड कल किराने की दुकान जाएगा,'' तो हम ऐसी
घटना के बारे में दावा कर रहे हैं जो अभी हुई नहीं है और जिसके सत्य-मूल्य पर
विवाद हो सकता है। वास्तव में, यह भी विवादास्पद है कि उस दावे को सबसे पहले
कोई सत्य-मूल्य \emph{दिया} भी जा सकता है या नहीं। यदि हम कठोर नियतिवादी हैं,
तो शायद हम यह मानने में सहज होंगे कि यह वाक्य वास्तव में सत्य या असत्य है,
उस घटना के घटित होने के निर्धारित समय से भी पहले---बस संभव है कि हमें अभी
उसका सत्य-मूल्य ज्ञात न हो। इसके विपरीत, हम भविष्य को सचमुच खुला मान सकते
हैं, जिसमें भावी आकस्मिक कथनों के सत्य-मूल्य अनिर्धारित रहते हैं।

वस्तुतः समय की संरचना और प्रकृति के बारे में ऐसी बहुत-सी प्रतिबद्धताएँ हमारे
कालिक प्रतिरूपों और तर्कशास्त्रीय ढाँचों के चुनाव में ही समाहित होती हैं।
उदाहरण के लिए, हम पूछ सकते हैं कि क्या हमें ऐसे प्रतिरूप बनाने चाहिए जिनमें
समय रैखिक, शाखित या चक्रीय भी हो। हमें यह भी तय करना पड़ सकता है कि हमारे
कालिक प्रतिरूपों में आरंभिक और अंतिम बिंदु होंगे या नहीं, और समय को पृथक
काल-क्षणों के द्वारा निरूपित किया जाएगा या एक सातत्य के रूप में।

\end{document}
